% Lösungen Einheit 3 von 3 – Skalarprodukt im Sachzusammenhang (zusammenbau v0.1)
\einheitenkopf{Einheit 3 von 3 · Skalarprodukt im Sachzusammenhang}

\erg{14}{a) Vektor \quad b) $(12 | 30 | 18)$; die zweite Komponente zählt die verkauften Cappuccinos. \quad c) Die Gesamtkosten des Einkaufs in Euro – eine Zahl, kein Vektor. \quad d) $\vec m = t \cdot (5 | 2 | 1)$, $t \cdot (20 + 12 + 8) = 168$, $t = 4{,}2$; $w = 21$, $g = 8{,}4$, $r = 4{,}2$}
\erg{15}{a) Fehler: die Produkte nicht addiert – das Skalarprodukt ist eine Zahl. Richtig: $8 + 6 + 7 = 21$ €. \quad b) Weil die Produkte der Komponenten addiert werden – übrig bleibt eine Summe, der Gesamtpreis. \quad c) $\vec m \circ \vec p = 28{,}80 + 30{,}00 + 24{,}00 = 82{,}80$ €; nein, 80 € reichen nicht.}
